% Transcription — fonds Alexandre Grothendieck, shelfmark 4, batch 1 (pages 1-11).
% Established from the facsimile archives/batches/4/batch-01.pdf (a mirror of
% https://grothendieck.umontpellier.fr/4.pdf), per the skill
% transcribe-grothendieck. The folder is eleven pages and this is its only
% batch. All eleven pages carry mathematics and his hand, and all are
% transcribed; none is skipped.
%
% The folder holds two things, as the inventory says. Page 1 is a single leaf
% of manuscript notes headed « 8.4.4. », in ink: it is the « feuille jointe
% (marquée 8.4.4) » that his marginal note on page 5 refers to, and the
% « en vertu de 8.4.4 » he writes on page 4. Pages 2 to 11 are copies of a
% typescript, §8 of a longer text, « Une estimation (heuristique) du nombre des
% valeurs propres de frobenius dans la cohomologie évanescente qui sont des
% unités p-adiques », which he went over in pencil and ink — insertions,
% renumberings, whole paragraphs boxed and cancelled, and marginal notes.
% Following folders 48 and 55, the numbered statements, the displayed formulas
% and all of his interventions are given, and the typed connective prose is
% given where his hand touches it and otherwise summarised in \note{} rather
% than recomposed. The typist's own overstrikes (runs of x over a mistyped
% word) are corrections of the typing and are not recorded; what he struck by
% hand is.
%
% The mathematics. For X smooth projective over a perfect field k of
% characteristic p, W = W(k), a crystalline cohomology H(X) is assumed with
% the universal-coefficient sequences (8.1.2) relating H^i(X) ⊗_W k,
% H^i_DR(X) and Tor_1^W(H^{i+1}(X), k). For a hyperplane or hypersurface
% section Y, the snake lemma gives (8.1.3), linking the evanescent
% cohomology E(Y) = Coker(H^{n-1}(X) → H^{n-1}(Y)) to its de Rham analogue.
% 8.2 (cancelled) dualises the torsion terms by Poincaré duality; 8.3 uses
% weak Lefschetz (admitted crystallinely) to make E(Y) torsion-free and the
% relevant maps independent of Y. 8.4 takes semi-simple parts for the
% p-linear Frobenius (exact functor), reads them as F_p-representations of
% Gal(k̄/k), and Lemme 8.4.3 shows H^i_DR(X) → H^i(X, O_X) is an isomorphism
% on semi-simple parts. « Théorème » 8.5 is the exact sequence (8.5.1)
% 0 → Ker v → (pTors^n(X))_ss → (E(Y) ⊗ k)_ss → E(Y, O_Y)_ss → 0, whence the
% dimension formula (8.5.5) and inequality (8.5.6): the number of Frobenius
% eigenvalues on the evanescent cohomology that are p-adic units is bounded
% below by the semi-simple rank of the coherent evanescent cohomology
% E(Y, O_Y). The remarks 8.5.7-8.5.11 treat injectivity of v, independence
% of Y over an open U_d, the local system on U_d, when φ is bijective, and
% Serre's Mexico example (φ not injective; X × Z gives φ not surjective for
% n ≥ 4). 8.6 asks whether the monodromy of the local system E(Y, O_Y)_ss is
% all of Aut(V), which with the « formule de Woodshole » and Čebotarev would
% control card Y(k') mod p, noting that nothing of the kind is proved even
% for plane curves. Page 1 records, for i > n and i < n, the nilpotence of F
% and V and the exact sequences (*) and (**) that page 4 and page 6 draw on,
% and fixes the notation: subscript ss for F, superscript ss for V.
%
% The typescript reads cleanly; the difficulty is the hand. The manuscript leaf
% (page 1) is a quick draft whose formulas carry. The pencil notes on the
% typescript are faint and fast, and several are left partly \ill{}: the
% slanted note at the foot of page 2, the faint note in the right margin of
% page 6 opening « Katz fait observer », the struck lines on page 8, the long
% interlinear note on page 11.
%
% Notation fixed once here. The structure sheaf, underlined on the typescript,
% is set \underline{O}; the semi-simple parts are set with an upright ss,
% subscript for F and superscript for V as page 1 defines them; the p-torsion
% « pTors » with a left subscript is set {}_{p}\mathrm{Tors}; « lib » (free
% part) is his.
%
% Readings flagged and not settled: the superscripts and the order of the
% terms in (*) on page 1; « Tor » versus « Tors » in line a) of page 1; the
% exponents i-n, i-n+1 in his marginal sequence on page 2; the hand label on
% the arrow w of (8.5.3), read u_{dξ}; « H^n » over a blotted exponent on
% page 8; the parenthesis « et, … si (pTors^n(X))_ss = 0 » on page 9; the note
% « Cor. 8.4.4 ter » in the margin of page 5, whose numbering is unsure; the
% placement of « H^1(Y, O_Y)_ss » on page 10.
%
% Pass: Opus 5 (claude-opus-5), 2026-09-14 — first pass, unchecked against the pages by a human.
\documentclass[11pt,a4paper]{article}
% Le préambule commun à toutes les transcriptions du fonds Grothendieck.
%
% Il tient en une idée : **l'incertitude se compose, elle ne se résout pas.**
% Une transcription qui lisse un mot douteux a détruit la seule chose pour
% laquelle on la faisait — un lecteur doit pouvoir savoir, à chaque endroit,
% ce qui était sur le feuillet et ce qui a été ajouté. Les six macros qui
% suivent sont donc l'appareil critique, et non de la décoration.
%
% Les mêmes commandes sont reconnues par `scripts/render.mjs`, qui produit la
% vue de lecture du site. Une macro ajoutée ici doit l'être aussi là-bas,
% faute de quoi le rendu échouera bruyamment — ce qui est voulu : un rendu
% silencieusement faux est pire qu'un rendu absent.

\ProvidesPackage{grothendieck}[2026/08/08 Transcriptions du fonds Grothendieck]

\usepackage{amsmath,amssymb,amsthm}
\usepackage{mathrsfs}
\usepackage{stmaryrd}
% Les diagrammes commutatifs sont partout dans ce fonds : c'est la langue
% même de la géométrie algébrique de Grothendieck, et une transcription qui
% les rendrait en prose ne transcrirait rien.
\usepackage{tikz-cd}
% Français par défaut — c'est la langue des feuillets — mais l'anglais est
% chargé aussi : la traduction et le résumé sont des documents anglais, et la
% typographie française (espaces avant les deux-points) y serait une faute.
% Ces documents basculent par \selectlanguage{english} en tête de corps.
\usepackage[english,french]{babel}
\usepackage{fontspec}
\usepackage{geometry}
\geometry{a4paper,margin=25mm,marginparwidth=28mm}
\usepackage{marginnote}
\usepackage{xcolor}
% ulem, not soul. `soul`'s \st and \ul are built on a character-by-character
% scan that XeLaTeX's font machinery breaks: the document fails with an
% undefined control sequence rather than degrading. ulem does the same two jobs
% and is Unicode-engine safe. `normalem` stops it hijacking \emph, which the
% transcriptions use for Grothendieck's own emphasis.
\usepackage[normalem]{ulem}
\usepackage{hyperref}
\hypersetup{colorlinks=true,linkcolor=black,urlcolor=black,pdfborder={0 0 0}}

% --- Métadonnées du lot -----------------------------------------------------
% Reprises telles quelles par le script de rendu, qui les affiche en tête de
% la vue de lecture. Elles fixent ce que le fichier prétend couvrir : sans
% elles, une transcription flotte sans point d'attache dans un fonds de seize
% mille feuillets.

\newcommand{\folder}[1]{\gdef\grothFolder{#1}}
\newcommand{\batch}[1]{\gdef\grothBatch{#1}}
\newcommand{\leaves}[2]{\gdef\grothFirst{#1}\gdef\grothLast{#2}}
\newcommand{\foldertitle}[1]{\gdef\grothFoldertitle{#1}}
\gdef\grothFolder{?}\gdef\grothBatch{?}\gdef\grothFirst{?}\gdef\grothLast{?}\gdef\grothFoldertitle{}

% --- Le résumé --------------------------------------------------------------
%
% Il ouvre la lecture modernisée, sous le titre « Résumé », et il est composé
% en petit. Ce n'est pas un ornement : c'est le seul passage du document qui
% ne soit pas des mathématiques, et un lecteur doit pouvoir le reconnaître
% avant de l'avoir lu — puis le sauter s'il connaît déjà le sujet, ou s'y
% arrêter s'il ne le connaît pas. Le filet qui le suit marque la reprise du
% corps.

\newenvironment{resume}
  {\small\setlength{\parskip}{0.45em}}
  {\par\medskip\hrule\medskip}

% --- Mots-clés ---------------------------------------------------------------
%
% En anglais, en clôture du résumé de la lecture modernisée : le vocabulaire
% moderne sous lequel chercher ce que les feuillets construisent. Le manifeste
% du site (`npm run manifest`) les extrait pour étiqueter la cote entière —
% cette ligne est la source unique des tags, il n'y a pas de fichier à côté.
\newcommand{\keywords}[1]{\par\smallskip\noindent
  {\footnotesize\textsc{Keywords} --- \textit{#1}}\par}

% --- L'appareil critique ----------------------------------------------------

% Le numéro de feuillet, en marge. C'est l'ancre qui permet de régler un
% désaccord sur une formule en regardant *un* feuillet plutôt qu'une chemise
% de six cents. Le site s'en sert aussi pour tourner les pages du fac-similé
% quand on fait défiler la transcription.
\newcommand{\leaf}[1]{%
  \marginnote{\footnotesize\color{gray!70}\textsf{#1}}%
  \label{leaf:#1}\ignorespaces}

% Illisible : on ne devine pas. Un mot inventé qui se lit comme les autres est
% le pire résultat possible.
\newcommand{\ill}{\textcolor{red!60!black}{[\,\ldots\,]}}

% Lecture douteuse : le mot est donné, et signalé comme incertain.
\newcommand{\uncertain}[1]{\uline{#1}}

% Ajout de l'éditeur — une lettre manquante, un mot sous-entendu.
\newcommand{\add}[1]{\textcolor{blue!50!black}{[#1]}}

% Biffé par Grothendieck lui-même. Ce qu'il a rayé fait partie du document :
% une rature dit souvent où la pensée a changé de direction.
\newcommand{\struck}[1]{\sout{#1}}

% Note du transcripteur, sur l'état matériel du feuillet.
\newcommand{\note}[1]{\par\smallskip{\small\color{gray!60!black}\textit{[#1]}}\par\smallskip}

% Note marginale de Grothendieck, distincte du corps du texte.
\newcommand{\marginal}[1]{\par\smallskip\noindent\hspace{1em}%
  {\small\color{gray!60!black}#1}\par\smallskip}

% --- Titre ------------------------------------------------------------------

\AtBeginDocument{%
  \begin{center}
    {\large\bfseries Fonds Alexandre Grothendieck --- cote n\textsuperscript{o}~\grothFolder}\\[2pt]
    {\itshape\grothFoldertitle}\\[4pt]
    {\small feuillets \grothFirst--\grothLast\ (lot \grothBatch)}\\[2pt]
    {\footnotesize\color{gray!60!black}Transcription établie d'après le fac-similé numérisé
     par l'Université de Montpellier.\\
     Les lectures incertaines sont soulignées, les illisibles notées [\,\ldots\,],
     les ajouts entre crochets.}
  \end{center}
  \vspace{1em}}

\endinput


\folder{4}
\batch{1}
\pages{1}{11}
\dating{[à partir de 1958]}
\watermark{Édition de démonstration}
\foldertitle{Torsion en cohomologie cristalline : notes manuscrites (s.d.), copies de tapuscrit annoté (s.d.)}

\begin{document}

\section*{8.4.4}

\note{notes manuscrites, à l'encre, sur un feuillet seul. C'est la « feuille
jointe (marquée 8.4.4) » à laquelle renvoie sa note marginale de la page 5, et
le « en vertu de 8.4.4 » de la page 4 ; elle est donnée ici à sa place dans le
dossier, avant le tapuscrit}

\page{1}

8.4.4. \quad \emph{NB}

a) $F$ opérant sur $H^{i}(X) \otimes_{W} k$ pour $i > n$ est \emph{nilpotent},
donc aussi opérant sur $\mathrm{Tors}^{i}(X) \otimes_{W} k$ ou sur
$\mathrm{Tor}^{W}_{1}(H^{i}(X), k) \simeq {}_{p}\mathrm{Tors}^{i}(X)$\note{la
lecture « Tor » et l'indice $W$ sont incertains : un premier signe est
repris à la main}. On a
\[
H^{n}(X, \underline{O}_{X})_{\mathrm{ss}} \simeq (H^{n}(X) \otimes_{W} k)_{\mathrm{ss}}
\]
d'où
\[
(*) \qquad 0 \to (\mathrm{Tors}^{n}(X) \otimes k)_{\mathrm{ss}}
\to H^{n}(X, \underline{O}_{X})_{\mathrm{ss}}
\to (H^{n}_{\mathrm{lib}}(X) \otimes k)_{\mathrm{ss}} \to 0
\]
\note{les exposants et l'ordre des termes de $(*)$ sont lus tels qu'ils se
présentent, sur une première tentative effacée à gauche ; ils ne sont pas
assurés}

b) $V$ opérant sur $H^{i}(X) \otimes_{W} k$ pour $i < n$ est nilpotent, et
aussi opérant sur $\mathrm{Tors}^{i}(X) \otimes_{W} k$ et
${}_{p}\mathrm{Tors}^{i}(X)$ pour $i \leq n$\note{« $i \leq n$ » est souligné
deux fois}. On a
\[
(**) \qquad 0 \to (H^{n}_{\mathrm{lib}}(X) \otimes k)^{\mathrm{ss}}
\to H^{0}(X, \Omega^{n}_{X})^{\mathrm{ss}}
\to {}_{p}\mathrm{Tors}^{n+1}(X)^{\mathrm{ss}} \to 0
\]

Sans doute $(*)$ et $(**)$ sont transposés l'un de l'autre \ill{}.

\emph{Notation} : l'indice $\mathrm{ss}$ indique qu'on prend la partie
semi-simple pour $F$, l'exposant $\mathrm{ss}$ \ill{} \add{pour $V$}.

\section*{8. Une estimation (heuristique) du nombre des valeurs propres de
frobenius dans la cohomologie évanescente qui sont des unités $p$-adiques}

\note{titre tapé en tête de la page 2, souligné à la main d'un bout à l'autre.
Les pages 2 à 11 sont des copies de ce tapuscrit, qu'il a reprises au crayon et
à l'encre}

\page{2}

\textbf{8.1.} Soit $X$ un schéma lisse et propre sur le corps parfait $k$
\add{de car.\ $p > 0$}\note{inséré à la main}, et soit $W$ l'anneau des
vecteurs de Witt sur $k$. Si on dispose d'une bonne théorie de la cohomologie
cristalline $H(X)$ (au sens des catégories dérivées), on devra avoir la relation
suivante avec la cohomologie de De Rham
\[
(8.1.1) \qquad H_{\mathrm{DR}}(X) = H(X) \otimes^{L}_{W} k
\]
(isomorphisme dans la catégorie dérivée des $k$-modules). On en déduit une
formule des coéfficients universels : suites exactes
\[
(8.1.2) \qquad 0 \to H^{i}(X) \otimes_{W} k \to H^{i}_{\mathrm{DR}}(X)
\to \mathrm{Tor}^{W}_{1}(H^{i+1}(X), k) \to 0 .
\]
Ce sont des suites exactes fonctorielles en $X$, tout comme l'isomorphisme
(8.1.1). \struck{Sauf erreur,} ceci marche dès maintenant pour $p \neq 2$,
grâce à Berthelot.\note{le « $L$ » de (8.1.1) est porté à la main au-dessus du
produit tensoriel}

\marginal{On conclut, grâce à 8.4.3,
$0 \to (H^{i}(X) \otimes_{W} k)_{\mathrm{ss}} \to
H^{i}(X, \underline{O}_{X})_{\mathrm{ss}}
\to ({}_{p}\mathrm{Tors}^{i+1}(X))_{\mathrm{ss}} \to 0$
(nul si $i > n$) \uncertain{(nul si $i \geq n$)} et
$0 \to (H^{i}(X) \otimes_{W} k)^{\mathrm{ss}} \to
H^{i-n}(X, \Omega^{n}_{X})^{\mathrm{ss}}
\to ({}_{p}\mathrm{Tors}^{i-n+1}(X))^{\mathrm{ss}} \to 0$,
où l'indice $\mathrm{ss}$ signifie la partie semi-simple pour
$F_{X} = F_{X}^{*}\sigma$, et où l'exposant $\mathrm{ss}$ \ill{}
$V_{X} = \sigma^{-1} F_{X*}$}\note{la première suite est écrite en interligne
après « Berthelot », la suite de la note monte dans la marge gauche, qu'il faut
tourner. Les exposants $i - n$ et $i - n + 1$ de la seconde suite sont
incertains ; « (nul si $i > n$) » est écrit sous le terme médian de la
première, « (nul si $i \geq n$) » au-dessus de son dernier terme}

Si maintenant $Y$ est une section hyperplane de \struck{degré assez élevé de}
$X$ (supposée maintenant $\subset \mathbf{P}^{r}_{k}$), on trouve une suite
exacte analogue pour $Y$, et un homomorphisme \add{$u = (u_{0}, u_{1}, u_{2})$}
de la première dans la seconde. Le diagramme du serpent nous donne alors une
suite exacte
\[
(8.1.3) \qquad \bigl(0 \to \mathrm{Ker}\, u_{0} \to \mathrm{Ker}\, u_{1} \to
\mathrm{Ker}\, u_{2} \to\bigr)\; E(Y) \otimes_{W} k \to
\underbrace{E_{\mathrm{DR}}(Y)}_{=\,\mathrm{Coker}\, u_{1}}
\to \mathrm{Coker}\, u_{2} \to 0 ,
\]
où on a posé
\[
E(Y) = \mathrm{Coker}\bigl(H^{n-1}(X) \xrightarrow{u_{0}} H^{n-1}(Y)\bigr),
\qquad
E_{\mathrm{DR}}(Y) = \mathrm{Coker}\bigl(H^{n-1}_{\mathrm{DR}}(X)
\xrightarrow{u_{1}} H^{n-1}_{\mathrm{DR}}(Y)\bigr),
\]
$X$ étant supposée connexe de dimension $n$, enfin $u_{2}$ désigne
l'application
\[
(8.1.4) \qquad u_{2} : \mathrm{Tor}^{W}_{1}(H^{n}(X), k) \to
\mathrm{Tor}^{W}_{1}(H^{n}(Y), k).
\]
\note{le triplet $u = (u_{0}, u_{1}, u_{2})$, les trois premiers termes de
(8.1.3) mis entre parenthèses, « $= \mathrm{Coker}\, u_{1}$ » et les étiquettes
des flèches sont de sa main}

\textbf{8.2.} Noter que $\mathrm{Tor}^{W}_{1}(-, k)$ est isomorphe à
$\mathrm{Hom}_{W}(k, -)$ (éléments de $-$ annulés par $pW$), et que
$\mathrm{Tors}^{n}(Y)$ (= sous-module de torsion de $H^{n}(Y)$) est dual (si on
admet la dualité de Poincaré habituelle pour la cohomologie cristalline, y
compris les formules pour la torsion) de $\mathrm{Tors}^{n-1}(Y)$, et\note{le
numéro 8.2 est porté à la main ; tout le paragraphe, jusqu'à la fin de son
premier alinéa en haut de la page 3, est barré de longues diagonales ; il est
donné en clair, la rature étant notée ici une fois}

\marginal{il est \uncertain{malcommode} d'utiliser ici la dualité de
Poincaré}\note{note écrite de biais dans l'angle inférieur gauche, en regard
du paragraphe 8.2 barré}

\page{3}

par suite le deuxième membre de (8.1.4) est canoniquement dual à
$\mathrm{Tors}^{n-1}(Y) \otimes_{W} k$. De même le premier membre est
canoniquement dual de $\mathrm{Tors}^{n+1}(X) \otimes_{W} k$, et enfin
l'homomorphisme $u_{2}$ est transposé de l'homomorphisme
\[
(8.2.1) \qquad u'_{2} : \mathrm{Tors}^{n-1}(Y) \otimes_{W} k \to
\mathrm{Tors}^{n+1}(X) \otimes_{W} k
\]
induit par l'homomorphisme de Gysin $H^{n-1}(Y) \to H^{n+1}(X)$. Par suite,
les termes $\mathrm{Ker}\, u_{2}$ et $\mathrm{Coker}\, u_{2}$ dans (8.1.3)
peuvent s'interpréter respectivement comme les duals de $\mathrm{Coker}\,
u'_{2}$ et $\mathrm{Ker}\, u'_{2}$.

\textbf{8.3.} Considérons l'homomorphisme induit par l'inclusion \add{$Y \to
X$},\quad $u_{0} : H^{n-1}(X) \to H^{n-1}(Y)$. \struck{Sauf erreur,} dans la
situation analogue \struck{sur le corps des complexes,} \add{en cohomologie
$\ell$-adique ($\ell \neq p$)}, \struck{les assertions} \add{le théorème} de
Lefschetz \struck{dans son livre « L'Analysis Situs \ldots{} »} \add{faible (pour
les coefficients $\mathbf{Z}/\ell^{n}\mathbf{Z}$)} impliquent que le conoyau
$E(Y)$ est sans torsion, \struck{ce} qui implique en particulier que $u_{0}$
induit un isomorphisme
\[
(8.3.1) \qquad \mathrm{Tors}^{n-1}(X) \simeq \mathrm{Tors}^{n-1}(Y).
\]
\note{les corrections de ce passage sont à la main ; la parenthèse sur les
coefficients est d'une lecture incertaine, et le « impliquent » tapé n'a pas été accordé}

Nous admettrons que ces relations restent valables dans le cas de la
cohomologie cristalline. Alors l'homomorphisme (8.2.1) s'interprète comme
l'homomorphisme
\[
(8.3.2) \qquad u''_{2} : \mathrm{Tors}^{n-1}(X) \otimes_{W} k \to
\mathrm{Tors}^{n+1}(X) \otimes_{W} k
\]
induit par l'homomorphisme $H^{n-1}(X) \to H^{n+1}(X)$, cup-produit par la
classe de polarisation de $X \subset \mathbf{P}^{r}_{k}$. La chose importante
pour nous est que ce morphisme (donc son Ker) ne dépend plus du choix de $Y$,
donc \struck{aussi du} \add{\ill{}} \add{le} dernier terme de (8.1.3) (qui est le
dual de ce Ker) ne dépend plus de $Y$.\note{au-dessus de « aussi du », biffé,
une insertion à la main elle-même en partie biffée, qui ne se lit pas}

\note{la suite de 8.3, tapée sans reprise, observe que $E(Y)$ étant sans
torsion, il a même rang que $E(Y) \otimes_{W} k$, de sorte que les valeurs
propres d'un endomorphisme de $E(Y)$ --- un frobenius --- se réduisent mod $p$
en celles de l'endomorphisme induit ; et, « petit luxe supplémentaire », que le
premier terme de (8.1.3) est nul}

\page{4}

\[
(8.3.3) \qquad \mathrm{Ker}\, u_{0} = 0 .
\]

\textbf{8.4.} \note{le tapuscrit introduit les applications semi-linéaires
Frob opérant fonctoriellement sur les termes de (8.1.3), vectoriels de
dimension finie sur $k$ à endomorphisme $p$-linéaire, et pour un tel $V$ la
partie « semi-simple » $V_{\mathrm{ss}}$, plus grand sous-espace stable sur
lequel $F_{V}$ est surjectif ($k$ parfait) ; il revient au même de dire que $F$
est surjectif sur ce sous-espace et \emph{nilpotent} sur le quotient --- ce mot
souligné à la main. La formation de $V_{\mathrm{ss}}$ est un foncteur exact en
$V$} Ainsi, la suite exacte (8.1.3) nous donne, compte tenu de (8.3.3), la suite
exacte d'espaces vectoriels à opérateurs $p$-linéaires semi-simples :
\[
(8.4.1) \qquad 0 \to (\mathrm{Ker}\, u_{1})_{\mathrm{ss}} \to
(\mathrm{Ker}\, u_{2})_{\mathrm{ss}} \to (E(Y) \otimes_{W} k)_{\mathrm{ss}}
\to E_{\mathrm{DR}}(Y)_{\mathrm{ss}} \to (\mathrm{Coker}\, u_{2})_{\mathrm{ss}}
\to 0 .
\]
\marginal{$(\mathrm{Ker}\, u_{2})_{\mathrm{ss}} =
({}_{p}\mathrm{Tors}^{n}(X))_{\mathrm{ss}}$ ; \quad
$(\mathrm{Coker}\, u_{2})_{\mathrm{ss}} = 0$ (en vertu de 8.4.4)}\note{la
première égalité est écrite au-dessus du terme, avec une flèche ; la seconde
est encerclée dans la marge droite, le « $0$ » relié au terme par un signe
d'égalité vertical}

\textbf{8.4.2.}\note{numéro porté à la main} D'après la théorie de structure des
vectoriels à opération $p$-linéaire semi-simple, ceux-ci peuvent s'interpréter
aussi comme des Modules libres de type fini sur le faisceau d'anneaux étale
constant $\underline{\mathbf{F}}_{p}$ (corps premier à $p$ éléments), ou encore
(une fois choisie une clôture séparable) comme des représentations de
$G = \mathrm{Gal}(\overline{k}/k)$ sur des vectoriels de dimension finie sur
$\mathbf{F}_{p}$. \note{le tapuscrit ajoute qu'il y a dans cette catégorie une
autodualité, qui s'interprète indifféremment dans l'un ou l'autre langage}

\page{5}

\note{alors les termes $(\mathrm{Coker}\, u_{2})_{\mathrm{ss}}$ et
$(\mathrm{Ker}\, u_{2})_{\mathrm{ss}}$ s'interprètent comme les duals de
$(\mathrm{Ker}\, u''_{2})_{\mathrm{ss}}$ et de $(\mathrm{Coker}\,
u''_{2})_{\mathrm{ss}}$, et $(\mathrm{Coker}\, u_{1})_{\mathrm{ss}}$,
$(\mathrm{Ker}\, u_{1})_{\mathrm{ss}}$ ont aussi une interprétation simple,
grâce au lemme qui suit}

\textbf{Lemme 8.4.3.} Soit $X$ un schéma propre sur un corps $k$, de car.\
$p > 0$. Alors pour tout $i \in \mathbf{Z}$, l'homomorphisme d'augmentation
\[
H^{i}_{\mathrm{DR}}(X) \to H^{i}(X, \underline{O}_{X})
\]
induit un isomorphisme sur les parties semi-simples (pour l'opération
$p$-linéaire de Frobenius).\note{le numéro tapé « 8.4.2 » est corrigé en
« 8.4.3 »}

Ceci résulte aussitôt de la suite spectrale de Hodge
\[
E_{1}^{a,b} = H^{b}(X, \underline{\Omega}^{a}_{X/k}) \Longrightarrow
H^{\bullet}_{\mathrm{DR}}(X),
\]
compte tenu que Frob opère $p$-linéairement sur cette suite spectrale et qu'il
est nul sur $E_{1}^{a,b}$ pour $a \neq 0$. Alors les
$(E_{1}^{0,b})_{\mathrm{ss}} = H^{b}(X, \underline{O}_{X})_{\mathrm{ss}}$ sont
des cycles universels pour les $d_{r}$, et on gagne \ldots{}

\marginal{Cor.\ 8.4.4 \uncertain{ter} (utilisé dans 8.1.\ill{}) les résultats
de la feuille jointe (marquée 8.4.4).}\note{écrit en montant dans la marge
gauche, en regard d'une grande accolade qui embrasse le lemme et sa
démonstration. La feuille jointe est la page 1}

Mettant ensemble les renseignements obtenus, \struck{et en même temps, pour des
raisons de commodité, transformant par dualité de Poincaré l'homomorphisme
$u''_{2}$ de (8.3.2), de sorte qu'on trouve (oubli dans 8.3 !)
(8.4.2) $u_{2}$ est can.\ isomorphe à} on trouve :\note{le passage biffé est
encadré et barré de traits obliques}

\textbf{« Théorème » 8.5.} En admettant les ingrédients précédents pour la
cohomologie cristalline (formule des coefficients universels, dualité de
Poincaré, théorème de Lefschetz \struck{« faible », y compris pour les groupes
de torsion} \add{sans négliger la torsion}), on obtient une suite exacte,
relative à une variété lisse projective connexe $X \subset \mathbf{P}^{r}_{k}$
de dimension $n$ sur un corps parfait de car.\ $p > 0$, et une section
hyperplane (ou par une hypersurface \add{de degré $d \geq 1$}, cela s'y
ramène \ldots{}) $Y$ lisse de dimension $n - 1$ :\note{« sans négliger la torsion »
et « de degré $d \geq 1$ » sont à la main ; au pied de la page, une première
frappe de la suite (8.5.1) est barrée et griffonnée}

\page{6}

\[
(8.5.1) \qquad 0 \to \mathrm{Ker}\, v \xrightarrow{w}
({}_{p}\mathrm{Tors}^{n}(X))_{\mathrm{ss}} \to
(E(Y) \otimes_{W} k)_{\mathrm{ss}} \xrightarrow{\varphi}
\underbrace{E(Y, \underline{O}_{Y})_{\mathrm{ss}}}_{\simeq\, \mathrm{Coker}\, v}
\to 0 ,
\]
\note{la suite est entièrement reprise à la main : le second terme tapé est
noirci et remplacé par $({}_{p}\mathrm{Tors}^{n}(X))_{\mathrm{ss}}$ encadré, les
étiquettes $w$ et $\varphi$ sont ajoutées, et la fin de la ligne tapée, après
$E(Y, \underline{O}_{Y})_{\mathrm{ss}}$, est noircie jusqu'au $0$}

où on désigne par $v$ l'homomorphisme de restriction
\[
(8.5.2) \qquad v : H^{n-1}(X, \underline{O}_{X})_{\mathrm{ss}} \to
H^{n-1}(Y, \underline{O}_{Y})_{\mathrm{ss}} ,
\]
(en particulier $\mathrm{Coker}\, v$ est la « cohomologie évanescente
cohérente » $E(Y, \underline{O}_{Y})$), et par $w$ l'homomorphisme
\[
(8.5.3) \qquad w : (\mathrm{Tors}^{n-1}(X) \otimes_{W} k)_{\mathrm{ss}}
\xrightarrow{u_{d\xi}} (\mathrm{Tors}^{n+1}(X) \otimes_{W} k)_{\mathrm{ss}}
\]
induit par l'homomorphisme $H^{n-1}(X) \to H^{n+1}(X)$, défini par le cup-produit
avec la classe de polarisation \add{$d\xi$} sur $X$.\note{le passage de « et par
$w$ » à « polarisation » est encadré et barré de traits obliques, et donné en
clair ;
l'étiquette $u_{d\xi}$ de la flèche est d'une lecture incertaine}

\marginal{Katz fait \uncertain{observer} \ill{} \ill{} \ill{} \ill{} \ill{}
sur la}\note{note au crayon très pâle dans la marge droite, en regard de
(8.5.3). Dans la marge gauche, un griffonnage appuyé et encerclé ne se lit pas}

Enfin, pour utiliser cette suite exacte (8.5.1), on tiendra compte du fait que
\[
(8.5.4) \qquad E(Y) \overset{\mathrm{dfn}}{=}
\mathrm{Coker}\bigl(H^{n-1}(X) \to H^{n-1}(Y)\bigr)
\]
est un $W$-module libre \emph{sans torsion} ; en particulier, si $k$ est un
corps fini, le rang de $(E(Y) \otimes_{W} k)_{\mathrm{ss}}$ est égal au nombre
des valeurs propres de Frobenius géométrique relatif à $k$ opérant sur $E^{n-1}(Y)$\note{le $E$ est frappé sur un $H$}
(la cohomologie cristalline évanescente de $Y$ ; ou ce qui revient au même, la
cohomologie $\ell$-adique pour n'importe quel $\ell \neq p$ --- sous réserve
de l'algébricité de l'opération $\Lambda$ sur les cohomologies cristallines et
$\ell$-adiques) \emph{qui sont des unités $p$-adiques} ; les valeurs propres en
question, réduites mod $p$, sont égales aux valeurs propres du frobenius
$p$-linéaire (relatif au corps de base fini $k$) opérant sur le terme
$E(Y) \otimes_{W} k$ de (8.5.1), ou si on préfère, en termes de
l'interprétation 8.4.2, aux valeurs propres de l'inverse du générateur de
frobenius de $G = \mathrm{Gal}(\overline{k}/k)$ opérant sur le vectoriel de
dimension finie sur $\mathbf{F}_{p}$ correspondant.

Notons d'autre part l'égalité suivante déduite de (8.5.1) :
\[
(8.5.5) \qquad \dim (E(Y) \otimes_{W} k)_{\mathrm{ss}} =
\dim E(Y, \underline{O}_{Y})_{\mathrm{ss}} +
\bigl(\dim ({}_{p}\mathrm{Tors}^{n}(X))_{\mathrm{ss}}
- \dim \mathrm{Ker}\, v\bigr) ,
\]
où le terme entre parenthèses du deuxième membre est $\geq 0$, de sorte
qu'on\note{après $E(Y, \underline{O}_{Y})_{\mathrm{ss}}$, un terme tapé est
noirci ; « $\mathrm{Coker}\, w$ » est biffé et remplacé à la main}

\page{7}

a aussi l'inégalité suivante (résultant déjà de la suite exacte des trois
derniers termes de (8.5.1)) :
\[
(8.5.6) \qquad \dim (E(Y) \otimes_{W} k)_{\mathrm{ss}} \geq
\dim E(Y, \underline{O}_{Y})_{\mathrm{ss}}
\]
$= \dim E(Y, \underline{O}_{Y})_{\mathrm{ss}} - t^{n-1}(X, p)$, où
$t^{i}(X, p)$ désigne le nombre de $p$-coéfficients de torsion de $X$ (défini
via la cohomologie cristalline de $X$ en dimension $i$).\note{le terme tapé
« $- \dim \mathrm{Ker}\, w$ » est noirci, et les deux lignes suivantes, de « $=$ » à « dimension
$i$ », sont encadrées et barrées ; le « $p$- » de « $p$-coéfficients » avait été ajouté à
la main avant la rature}

\textbf{Remarques 8.5.7.} Je pense que l'homomorphisme $v$ de (8.5.2) est
injectif (ce qui simplifie en conséquence (8.5.1) et (8.5.5), puisque on
aurait $\mathrm{Ker}\, v = 0$) ; c'est vrai en tous cas si \add{$n = 2$, ou si}
$Y$ est une section par une hypersurface de degré $d$ assez grand pour qu'on
ait
\[
H^{n-1}(X, \underline{O}_{X}(-d)) = 0 .
\]

D'autre part, on peut si on veut dans 8.5.5 bloquer ensemble les deux
termes $\dim \mathrm{Coker}\, w - \dim \mathrm{Ker}\, w$, pour trouver la
différence des rangs des deux membres de (8.5.3).\note{alinéa encadré et barré
de boucles, donné en clair}

\textbf{8.5.8.} Lorsque l'on applique (8.5.5) en prenant pour $Y$ une section
de $X$ avec une hypersurface de degré $d$ « assez générale » (de telle façon
que le rang semi-simple de $E(Y, \underline{O}_{Y})$ ait la valeur maxima),
alors on trouve qu'au terme $- \dim \mathrm{Ker}\, v$ près (sans doute toujours
nul, cf.\ 8.5.7) le deuxième membre est \emph{indépendant du choix particulier
de $Y$}. Cette remarque s'applique en particulier (8.5.7) si $d$ est assez
grand. On en conclut que si $k$ est fini, on peut trouver un ouvert $U_{d}$ de
l'espace projectif des hypersurfaces de degré $d$, tel que pour un point
variable de $U_{d}$ à coéfficients dans une extension finie $k'$ de $k$, le
nombre de valeurs propres de Frobenius opérant sur la cohomologie évanescente de
$Y_{\overline{k}}$ et qui sont des unités $p$-adiques soit constant. Nous
désirons mon-

\page{8}

trer que pour $d$ assez grand, ce nombre est non nul (ce qui impliquera alors
--- et est sans doute équivalent --- que le motif $E(Y)$ pour tous ces $Y$ est
de \emph{niveau} maximum $n - 1$, ou encore que le motif $E(Y)$ pour $Y$
générique est de niveau $n - 1$). Sous la forme de (8.5.6), on voit que ceci
est assuré dès que l'on sait que
\[
\dim E(Y, \underline{O}_{Y})_{\mathrm{ss}} > 0 \qquad \text{pour} \quad
d \to +\infty ,
\]
où le dim est pris bien entendu avec sa valeur générique.\note{le « $> 0$ »
est à la main, sur un terme tapé noirci}
On sait seulement (cf.\ n° 7) que $\dim E(Y,
\underline{O}_{Y})_{\mathrm{ss}} \to$ \ill{} pour $d$ grand, et que
$\varphi \neq 0$ pour $d$ grand (donc $E(Y)$ de niveau $n - 1$) dès que l'on
sait que $\dim ({}_{p}\mathrm{Tors}^{n+1}(X) \otimes k)_{\mathrm{ss}} <$
\uncertain{$p$} \ill{}\note{deux lignes et demie à la main, serrées sous la
ligne tapée et biffées, données en clair ; la lecture en est incertaine d'un bout à l'autre}

\marginal{or cela est établi au n° 7}\note{dans la marge gauche, écrit de biais,
avec une flèche vers la phrase biffée}

\textbf{8.5.9.} Le corps de base $k$ étant à nouveau quelconque (même pas
nécessairement parfait), désignons, pour un degré $d$ donné, par $U_{d}$
l'ouvert de l'espace projectif des hypersurfaces de degré $d$ dans
$\mathbf{P}^{r}_{k}$, formé des hypersurfaces dont l'intersection avec $X$ est
lisse de dimension $n - 1$, et telles que $E(Y, \underline{O}_{Y})_{\mathrm{ss}}$
soit de rang maximum. Alors les suites exactes (8.5.1), pour les $Y$ provenant
de points de $U_{d}$, proviennent d'une suite exacte de Modules localement
libres sur $U_{d}$ à opération $p$-linéaire semi-simple, ou de façon
équivalente (cf 8.4.2) à une suite exacte de systèmes locaux de vectoriels sur
$\mathbf{F}_{p}$ sur $U_{d}$ pour la topologie étale, ou enfin, une fois choisi
un point géométrique de $U_{d}$ (par exemple via une clôture séparable du point
générique) une suite exacte de représentations de $G = \pi_{1}(U_{d})$ dans des
vectoriels de dimension finie sur $\mathbf{F}_{p}$. On trouve un homomorphisme
\add{$\varphi$} du système local défini par les $E(Y) \otimes_{W} k'$
\struck{dans} \add{sur} celui défini par les $E(Y, \underline{O}_{Y})$, dont le
noyau \struck{et le conoyau sont des} \add{est un} système local
\emph{constant} (i.e.\ provenant du corps de base $k$). Lorsque $k$ est fini, on
peut supposer, quitte à passer à une extension finie de $k$, que les valeurs
propres du frobenius relatif à $k$ opérant sur ce dernier sont toutes égales à
$1$ (il suffit que les valeurs propres de frobenius relatif à $k$ opérant sur
$H^{n}(X, \underline{O}_{X})_{\mathrm{ss}}$ \struck{et sur les deux membres de
(8.5.3)} soient\note{l'exposant de $H^{n}$ est repris à la main sur un premier
chiffre noirci ; \uncertain{$n$}}

\page{9}

toutes égales à $1$). On trouve alors une description complète de la variation
des valeurs propres mod $p$ de frobenius opérant sur $E(Y)$, pour $Y$ variable
avec un point de $U_{d}$, en termes essentiellement de la seule représentation
de $G$ dans un vectoriel de dimension finie sur $\mathbf{F}_{p}$, correspondant
aux $E(Y, \underline{O}_{Y})_{\mathrm{ss}}$ (donc essentiellement aux « matrices
de Hasse-Witt des $Y$ en dimension critique), --- et bien entendu de la
distribution des éléments de frobenius dans $G$ en termes d'un point variable de
$U_{d}$.

\textbf{8.5.10.} L'homomorphisme fondamental $\varphi$ de (8.5.1) est
\struck{surjectif (resp.} injectif \add{donc bijectif} (si et seulement si $w$
dans 8.5.1 est \struck{injectif (resp.\ si $w$ est} surjectif, \add{et, \ill{}
si $({}_{p}\mathrm{Tors}^{n}(X))_{\mathrm{ss}} = 0$,} le seulement si étant vrai
sous la réserve $\mathrm{Ker}\, v = 0$ de 8.5.7) ; il suffit en particulier pour
ceci que $X$ n'ait pas de $p$-torsion en dimension $n$ \struck{(resp.\ en
dimension $n + 1$} (ou ce qui revient au même par dualité, en dimension
\uncertain{$n + 1$})\note{les « resp. » sont biffés à la main tout au long de
l'alinéa, ramenant l'énoncé au seul cas injectif ; l'insertion « et, \ldots{} si
$({}_{p}\mathrm{Tors}^{n}(X))_{\mathrm{ss}} = 0$ » est écrite en interligne, et
l'exposant de la dernière parenthèse est repris sur un chiffre noirci}.
Lorsque $X$ se relève (avec sa polarisation) en caractéristique nulle, cette
condition suffisante doit aussi pouvoir s'exprimer (une fois qu'on disposera de
fondements satisfaisants pour la cohomologie cristalline) par la condition
analogue sur \add{la torsion de} la cohomologie $p$-\emph{adique} de la variété
relevée (sur un corps algébriquement clos, ou si on veut, sur le corps des
complexes \ldots{}) : on réécrit (8.5.3) en termes de $p$-torsion de cette
variété tensorisée par $\mathbf{Z}/p\mathbf{Z}$ (mais en supprimant les indices
ss, qui n'ont plus de sens), et on demande que $w_{p}$ ainsi défini soit
surjectif resp.\ injectif \ldots{}\note{les accords au singulier
(« cette condition suffisante doit », « la condition analogue ») sont faits à
la main}

Je pense que même si $n = \dim X = 2$, il peut arriver que $\varphi$ ne
soit pas surjectif, et qu'on ne peut pas dans l'inégalité (8.5.6) supprimer le
terme correctif $- \dim \mathrm{Ker}\, w$ ou $t^{n-1}(X, p)$. En effet, si $X$
est une\note{de « on réécrit » au bas de la page, et le haut de la page 10
jusqu'à la condition c), tout est barré de grands zigzags ; le passage est
donné en clair}

\page{10}

Noter d'autre part qu'on a équivalence entre les conditions suivantes :
\begin{enumerate}
\item[a)] L'homomorphisme $\varphi$ est surjectif (resp.\ injectif) pour tout
$d$ et tout point $s \in U_{d}$ (définissant un $Y$ correspondant).
\item[b)] Même condition, mais avec \emph{un} $d$ fixé multiple de $p$, et tel
que l'on ait $\mathrm{Ker}\, v = 0$ (condition sans doute automatiquement
vérifiée (8.5.7)).
\item[c)] $(\mathrm{Tors}^{n-1}(X) \otimes_{W} k)_{\mathrm{ss}} = 0$ (resp.\
$(\mathrm{Tors}^{n+1}(X) \otimes_{W} k)_{\mathrm{ss}} = 0$).
\end{enumerate}

\textbf{8.5.11.}\note{numéro porté à la main} \struck{Il} serait intéressant de
donner des exemples où \struck{ces conditions ne sont pas satisfaites. Notez
que la condition non respée est toujours satisfaite si $n = 2$ ($X$ une
surface) ? Dans le cas général,} \add{$\varphi$ n'est pas bijectif} on a
l'inégalité
\[
\begin{aligned}
(8.5.11.1) \qquad \dim_{k} H^{n-1}(X, \underline{O})_{\mathrm{ss}} &=
\dim_{\mathbf{F}_{p}} H^{n-1}(\overline{X}, \mathbf{Z}/p\mathbf{Z}) \\
&\geq
\dim\bigl((H^{n-1}(X)/\mathrm{Tors}^{n-1}(X)) \otimes k\bigr)_{\mathrm{ss}}
\end{aligned}
\]
(grâce à 8.1.2 et 8.4.3 ; en fait cela marche en toute dimension $i$, pas
seulement $i = n - 1$), et dire que $\varphi$ est bijectif signifie que
l'inégalité précédente est une égalité. L'exemple de Serre au symposium de
Mexico fournit un $X$ (quotient d'une intersection complète par un groupe
$G = \mathbf{Z}/p\mathbf{Z}$ opérant sans points fixes) pour lequel on a
$H^{1}(X) = 0$, mais $H^{1}(X, \mathbf{Z}/p\mathbf{Z}) \neq \mathbf{Z}/{}$\add{$p\mathbf{Z}$}\note{le bord du feuillet coupe la fin de la
formule}. Quitte à couper avec des hyperplans, on trouve le même exemple avec
$\dim X = 2$. Pour cet exemple, d'après 8.1.2, on aura
\[
({}_{p}\mathrm{Tors}^{2}(X))_{\mathrm{ss}} \simeq
H^{1}(X, \mathbf{Z}/p\mathbf{Z}) \otimes k = \mathrm{Ker}\, \varphi
\]
(quel que soit $d$ et la section $Y$ de $X$ de degré $d$, permettant de définir
$\varphi$), qui est non nul, donc $\varphi$ n'est pas injectif.\note{l'indice
$p$ de ${}_{p}\mathrm{Tors}$, l'exposant $2$ repris sur un autre chiffre et
« $\mathrm{Ker}\, \varphi$ » sont à la main ; au-dessus, il écrit
« $H^{1}(Y, \underline{O}_{Y})_{\mathrm{ss}}$ » avec un signe d'égalité
vertical dont le terme d'attache n'est pas assuré}

Si maintenant $Z$ est un schéma lisse projectif sur $k$ de dimension
$n - 2$ ($n$ donné $\geq 4$), à cohomologie cristalline sans torsion et
$(H^{n-4}(Z) \otimes k)_{\mathrm{ss}} \neq 0$ (par exemple
$\mathbf{P}^{n-2}_{k}$ si $n$ pair \ldots{}), alors, posant $X' = X \times Z$, on
voit (sous réserve de la \emph{formule de Kunneth} en cohomologie cristalline)
que $(\mathrm{Tors}^{n-1}(X') \otimes k)_{\mathrm{ss}}$ est non nul, ce qui
fournit des exemples où $\varphi$ n'est pas surjectif, pour tout $n \geq 4$. Il
y a sans doute aussi des exemples avec $n = 3$, en cherchant bien \ldots{}\note{le
paragraphe est encadré et barré de diagonales, et donné en clair ; le « $4$ » y est repris à la
main sur un « $3$ »}

\page{11}

\textbf{8.6.} La suite exacte (8.5.9) de systèmes locaux sur $U_{d}$
correspondant aux suites exactes (8.5.1) met en évidence que le système local
correspondant aux $(E(Y) \otimes_{W} k)_{\mathrm{ss}}$ \struck{et aux
$E(Y, \underline{O}_{Y})_{\mathrm{ss}}$ ont} \add{a}, en plus d'un sous-système
\struck{resp.\ d'un système quotient} \add{du premier} de nature essentiellement
triviale (car provenant du corps de base, donc constant du point de vue
géométrique) \uncertain{un} composa\supplied{nt} \struck{commun $\mathrm{Coim}\,
\varphi = \mathrm{Im}\, \varphi$} \add{$E(Y, \underline{O}_{Y})_{\mathrm{ss}}$},
dont on s'attend qu'il donne lieu à une représentation tout à fait non triviale
de $G$ (NB si on n'aime pas les groupes fondamentaux généraux et les systèmes
locaux, on peut tout interpréter en termes de la suite exacte 8.5.1
correspondant au \emph{point générique} $s$ de l'espace des hypersurfaces de
degré $d$ dans $\mathbf{P}^{r}_{k}$, et remplacer $G$ par le groupe
$\mathrm{Gal}(\overline{k(s)}/k(s))$, --- mais on perd alors les frobenius
\ldots{}). On aimerait pouvoir préciser ce point. Par exemple, \add{(si $n \geq 2$,
cette représentation est-elle \emph{irréductible}, ou même} l'image du groupe
fondamental \add{$G$} dans $\mathrm{Aut}(V)$ ($V$ étant le vectoriel de
dimension finie sur $\mathbf{F}_{p}$ où s'effectue la représentation envisagée
de $G$) est-elle égale à $\mathrm{Aut}(V)$ tout entier ?\add{)}
\struck{Si oui,} alors dès que \add{l'on a} $\dim V \geq 2$ (ce qui est sans
doute vérifié dès que $d$ est assez grand, $\dim V$ tendant vers $\infty$ avec
$d$, cf.\ 8.5.8), tout élément de $\mathbf{F}_{p}$ est trace d'un élément de
$G$ opérant sur $V$.\note{les insertions sont au crayon ; « irréductible » est
souligné}

\marginal{(Ce n'est pas \struck{\ill{}} \ill{} lorsque
$X = \mathbf{P}^{r-1}_{k}$, cas qu'on peut traiter trivialement.) Lorsque la
condition \ill{}}\note{longue note au crayon en interligne au-dessus de
« entier ? », qui déborde vers la droite ; en dessous, une reprise en plus petit
ne se lit pas}

De ceci, de la « formule de Woodshole » \struck{\ill{}} et du théorème de
Čebotareff, on déduirait alors, si le corps $k$ est fini, qu'on peut, pour
n'importe quel ouvert non vide $U$ de l'espace des hypersurfaces de degré $d$
dans $\mathbf{P}^{r}_{k}$ ($d$ étant choisi comme dessus) trouver un point $s$
de $U$ à coéfficients dans une extension finie $k'$ de $k$, définissant une
section $Y$ telle que $\mathrm{card}\, Y(k')$ soit congru mod $p$ à un élément
de $\mathbf{F}_{p}$ donné à l'avance. Au moment d'écrire ces lignes, et malgré
l'assistance de gens compétents comme Katz, il ne semble pas qu'un résultat de
cette nature soit prouvé, même si $X = \mathbf{P}^{2}$ \supplied{(}i.e.\ pour des courbes planes \add{$Y$} de grand degré\supplied{)}.
Ce résultat serait déjà bien intéressant, car il impliquerait déjà une
réponse affirmative à la conjecture du n° 6 sur le niveau de la cohomologie
évanescente (dont il est question aussi dans 8.5.8).\note{« de la formule de
Woodshole » est tapé en interligne, suivi de mots biffés à la main. Au-dessus de
$X = \mathbf{P}^{2}$ il écrit ce qui semble être « $2 = r$ », et les deux dernières
lignes et demie, « Ce résultat \ldots{} 8.5.8) », sont encadrées et barrées de
diagonales}

\end{document}
